\documentclass[11pt]{article}
\usepackage{amsmath}
\usepackage{geometry}
\geometry{a4paper, margin=1in}
\usepackage{enumitem}
\usepackage{booktabs}

\title{GED Math Worksheet: Tables, Graphs, and Interpreting Data - Solutions [cite: 1]}
\author{}
\date{}

\begin{document}

\maketitle

\noindent\textbf{Name:} \rule{5cm}{0.4pt} [cite: 2]\\
\textbf{Date:} \rule{5cm}{0.4pt} [cite: 3]\\

\noindent\textbf{Instructions:} Answer all 30 questions. Show your work where applicable. You may use a calculator for these problems. [cite: 4]

\vspace{0.5cm}

\section*{Part A: Tables (Questions 1-6) [cite: 5]}

\noindent \textbf{Questions 1-3 refer to the following table: [cite: 6]}

\begin{table}[h]
\centering
\begin{tabular}{lcccc}
\toprule
\textbf{Month} & \textbf{Jan} & \textbf{Feb} & \textbf{Mar} & \textbf{Apr} \\ \midrule
Sales (\$1000s) & 45 & 52 & 38 & 61 \\ \bottomrule
\end{tabular}
\end{table}
\vspace{0.2cm} [cite: 7]

\begin{enumerate}
    \item \textbf{What were the total sales for the four months? [cite: 8]} \\
    \textbf{Solution:} Total Sales = $45 + 52 + 38 + 61 = 196$. \\
    Since the values are in \$1000s, the total is \textbf{\$196,000}.
    
    \item \textbf{What is the average monthly sales for this period? [cite: 9]} \\
    \textbf{Solution:} Average = $\frac{196}{4} = 49$. \\
    The average monthly sales is \textbf{\$49,000}.
    
    \item \textbf{By how much did sales increase from February to April? [cite: 10]} \\
    \textbf{Solution:} April Sales (\$61,000) - February Sales (\$52,000) = 9,000. \\
    Sales increased by \textbf{\$9,000}.
\end{enumerate}

\vspace{0.5cm}
\noindent \textbf{Questions 4-6 refer to the following table: [cite: 11]}

\begin{table}[h]
\centering
\begin{tabular}{lccc}
\toprule
\textbf{Product} & \textbf{Price} & \textbf{Quantity Sold} & \textbf{Revenue} \\ \midrule
Widget A & \$8 & 125 & \$1,000 \\
Widget B & \$12 & 87 & \$1,044 \\
Widget C & \$15 & 92 & \$1,380 \\
Widget D & \$10 & 156 & \$1,560 \\ \bottomrule
\end{tabular}
\end{table}
\vspace{0.2cm} [cite: 13]

\begin{enumerate}[resume]
    \item \textbf{Which product generated the highest revenue? [cite: 14]} \\
    \textbf{Solution:} Comparing revenues: Widget A (\$1,000), Widget B (\$1,044), Widget C (\$1,380), Widget D (\$1,560). \\
    \textbf{Widget D} generated the highest revenue.
    
    \item \textbf{If Widget B's price increases by 20\%, what is the new price? [cite: 15]} \\
    \textbf{Solution:} Current price of Widget B = \$12. \\
    Increase = $\$12 \times 0.20 = \$2.40$. \\
    New Price = $\$12 + \$2.40 = \textbf{\$14.40}$.
    
    \item \textbf{What is the total number of widgets sold across all products? [cite: 16]} \\
    \textbf{Solution:} Total = $125 + 87 + 92 + 156 = \textbf{460}$ widgets.
\end{enumerate}

\section*{Part B: Bar Graphs (Questions 7-12) [cite: 17]}

\noindent \textbf{Questions 7-9 refer to the following data about monthly temperatures (in $^\circ$F): [cite: 18]} \\
Month: Jan (32), Feb (35), Mar (45), Apr (58), May (72), Jun (85) [cite: 18]

\begin{enumerate}[resume]
    \item \textbf{What is the difference in temperature between the warmest and coldest month listed? [cite: 19]} \\
    \textbf{Solution:} Warmest (June) = $85^\circ$F, Coldest (January) = $32^\circ$F. \\
    Difference = 85 - 32 = \textbf{$53^\circ F$}.
    
    \item \textbf{What is the average temperature for the six months shown? [cite: 20]} \\
    \textbf{Solution:} Average = $\frac{32 + 35 + 45 + 58 + 72 + 85}{6} = \frac{327}{6} = 54.5^\circ F$.
    
    \item \textbf{Between which two consecutive months is the greatest temperature increase? [cite: 21]} \\
    \textbf{Solution:} \\
    Jan to Feb: $+3^\circ$F \\
    Feb to Mar: $+10^\circ$F \\
    Mar to Apr: $+13^\circ$F \\
    Apr to May: $+14^\circ$F \\
    May to Jun: $+13^\circ$F \\
    The greatest increase ($+14^\circ$F) is between \textbf{April and May}.
\end{enumerate}

\vspace{0.5cm}
\noindent \textbf{Questions 10-12 refer to the following data about student enrollment in classes: [cite: 23]} \\
Math (120), Science (95), English (110), History (85), Art (60) [cite: 23]

\begin{enumerate}[resume]
    \item \textbf{How many more students are enrolled in Math than in Art? [cite: 24]} \\
    \textbf{Solution:} Math (120) - Art (60) = \textbf{60} students.
    
    \item \textbf{What is the total enrollment across all five classes? [cite: 25]} \\
    \textbf{Solution:} Total = $120 + 95 + 110 + 85 + 60 = \textbf{470}$ students.
    
    \item \textbf{What percentage of total enrollment is in Science? (Round to the nearest whole percent) [cite: 26]} \\
    \textbf{Solution:} Percentage = $\frac{95}{470} \approx 0.2021$. \\
    Rounded to the nearest whole percent, it is \textbf{20\%}.
\end{enumerate}

\section*{Part C: Line Graphs (Questions 13-17) [cite: 27]}

\noindent \textbf{Questions 13-17 refer to the following data showing website traffic (in thousands of visitors) over six weeks: [cite: 28]} \\
Week 1 (12), Week 2 (15), Week 3 (14), Week 4 (18), Week 5 (22), Week 6 (20) [cite: 29]

\begin{enumerate}[resume]
    \item \textbf{During which week did the website experience the highest traffic? [cite: 30]} \\
    \textbf{Solution:} The highest traffic occurred in \textbf{Week 5} (22,000 visitors).
    
    \item \textbf{What was the change in visitors from Week 1 to Week 6? [cite: 31]} \\
    \textbf{Solution:} Week 6 (20) - Week 1 (12) = 8. \\
    There was an \textbf{8,000 visitor increase}.
    
    \item \textbf{During which consecutive weeks was there a decrease in visitors? [cite: 32]} \\
    \textbf{Solution:} Traffic decreased from \textbf{Week 2 to Week 3} (15k to 14k) and from \textbf{Week 5 to Week 6} (22k to 20k).
    
    \item \textbf{What is the average number of visitors across all six weeks? [cite: 34]} \\
    \textbf{Solution:} Average = $\frac{12 + 15 + 14 + 18 + 22 + 20}{6} = \frac{101}{6} \approx 16.833$. \\
    The average is approximately \textbf{16,833 visitors}.
    
    \item \textbf{If the trend from Week 4 to Week 5 continues, how many visitors would you predict for Week 7? [cite: 35]} \\
    \textbf{Solution:} The trend from Week 4 to Week 5 is an increase of 4 (from 18 to 22). \\
    If we apply an increase of +4 to the Week 6 data point (20), the prediction for Week 7 would be $20 + 4 = \textbf{24}$ (or 24,000 visitors).
\end{enumerate}

\section*{Part D: Circle Graphs (Questions 18-22) [cite: 36]}

\noindent \textbf{Questions 18-22 refer to the following data about a budget breakdown for a household (\$3,600 dollars total): [cite: 37]} \\
Housing (40\%), Food (25\%), Transportation (18\%), Utilities (12\%), Entertainment (5\%) [cite: 37]

\begin{enumerate}[resume]
    \item \textbf{How much money is allocated to Housing? [cite: 38]} \\
    \textbf{Solution:} Housing = $\$3,600 \times 0.40 = \textbf{\$1,440}$.
    
    \item \textbf{What is the difference between the amount spent on Food and Entertainment? [cite: 39]} \\
    \textbf{Solution:} \\
    Food = $\$3,600 \times 0.25 = \$900$. \\
    Entertainment = $\$3,600 \times 0.05 = \$180$. \\
    Difference = $\$900 - \$180 = \textbf{\$720}$.
    
    \item \textbf{If the Entertainment budget is increased by 3\%, how much money would be allocated to it? [cite: 40]} \\
    \textbf{Solution:} Assuming the budget percentage increases by 3 percentage points (to 8\%): \\
    New allocation = $\$3,600 \times 0.08 = \textbf{\$288}$. \\
    *(Note: If interpreted as a 3\% increase on the \$180 amount, it would be $\$180 \times 1.03 = \$185.40$, but standard GED phrasing implies adjusting the percent share).*
    
    \item \textbf{What percentage of the budget is spent on Housing and Transportation combined? [cite: 41]} \\
    \textbf{Solution:} $40\% \text{ (Housing)} + 18\% \text{ (Transportation)} = \textbf{58\%}$.
    
    \item \textbf{If total household income increases to \$4,000 but percentages remain the same, how much will be spent on Food? [cite: 42]} \\
    \textbf{Solution:} New Food Budget = $\$4,000 \times 0.25 = \textbf{\$1,000}$.
\end{enumerate}

\section*{Part E: Scatter Plots (Questions 23-26) [cite: 44]}

\noindent \textbf{Questions 23-26 refer to the following data showing the relationship between study hours and test scores: [cite: 45]} \\
Study Hours: 1, 2, 3, 4, 5, 6 \\ Test Scores: 62, 68, 75, 81, 87, 92 [cite: 46]

\begin{enumerate}[resume]
    \item \textbf{What type of correlation does this data show? [cite: 47]} \\
    \textbf{Solution:} As study hours increase, test scores consistently increase. This shows a \textbf{strong positive correlation}.
    
    \item \textbf{Based on the pattern, what would you estimate the test score to be for 7 hours of study? [cite: 48]} \\
    \textbf{Solution:} The scores increase by roughly 5 to 7 points per hour. \\
    Estimating from the 6-hour mark (92) with an average increase of $+6$: $92 + 6 = \textbf{98}$.
    
    \item \textbf{What is the range of test scores shown in the data? [cite: 49]} \\
    \textbf{Solution:} Highest (92) - Lowest (62) = \textbf{30}.
    
    \item \textbf{Is there an outlier in this dataset? Explain your answer. [cite: 50]} \\
    \textbf{Solution:} \textbf{No}, there are no outliers. All test score increases are relatively consistent and follow a steady upward trend without any data points deviating significantly.
\end{enumerate}

\section*{Part F: Histograms (Questions 27-29) [cite: 51]}

\noindent \textbf{Questions 27-29 refer to the following histogram data showing distribution of ages at a community center: [cite: 52]} \\
Age 18-25 (8), Age 26-35 (15), Age 36-45 (12), Age 46-55 (10), Age 56-65 (5) [cite: 53]

\begin{enumerate}[resume]
    \item \textbf{What is the total number of people represented in the histogram? [cite: 54]} \\
    \textbf{Solution:} Total = $8 + 15 + 12 + 10 + 5 = \textbf{50}$ people.
    
    \item \textbf{Which age group has the highest frequency? [cite: 55]} \\
    \textbf{Solution:} The highest frequency is 15, which is the \textbf{26-35} age group.
    
    \item \textbf{What percentage of people are age 46 or older? (Round to the nearest whole percent) [cite: 57]} \\
    \textbf{Solution:} People 46 or older = $10 \text{ (from 46-55)} + 5 \text{ (from 56-65)} = 15$. \\
    Percentage = $\frac{15}{50} = 0.30$. \\
    The percentage is \textbf{30\%}.
\end{enumerate}

\section*{Part G: Box Plots (Question 30) [cite: 58]}

\noindent \textbf{Question 30 refers to the following dataset: 12, 15, 18, 22, 25, 28, 30, 32, 35 [cite: 59]}

\begin{enumerate}[resume]
    \item \textbf{Find the median, lower quartile, upper quartile, and interquartile range for this dataset. [cite: 59]} \\
    \textbf{Solution:} The dataset has 9 ordered numbers.
    \begin{itemize}
        \item \textbf{Median (Q2):} The middle number is \textbf{25}.
        \item \textbf{Lower Quartile (Q1):} The median of the lower half (12, 15, 18, 22) is $\frac{15 + 18}{2} = \textbf{16.5}$.
        \item \textbf{Upper Quartile (Q3):} The median of the upper half (28, 30, 32, 35) is $\frac{30 + 32}{2} = \textbf{31}$.
        \item \textbf{Interquartile Range (IQR):} $Q3 - Q1 = 31 - 16.5 = \textbf{14.5}$.
    \end{itemize}
\end{enumerate}

\vspace{1cm}
\begin{center}
    \textbf{End of Worksheet [cite: 60]}
\end{center}

\end{document}